\documentclass[10pt,a4paper]{amsart}
\usepackage[margin=19mm]{geometry}
\usepackage{amsmath,amssymb,mathtools}
\usepackage[scr=rsfs,cal=euler]{mathalfa}
\usepackage{xfrac}
\usepackage[unicode,hidelinks,bookmarks=false]{hyperref}
\newcommand{\ie}{i.e.\ }
\newcommand{\cf}{cf.\ }
\newcommand{\resp}{respectively\ }
\newcommand{\Subsection}{\subsection}
\providecommand{\nobreakdash}{\nobreak\hbox{-}}
\setlength{\emergencystretch}{2em}
\multlinegap0pt
\usepackage{bm}
\usepackage{bbm}
\usepackage{dsfont}
\usepackage{xspace}
\usepackage{cancel}
\usepackage{mathtools}
\usepackage{amsthm}
\makeatletter
\@ifundefined{thm}{\newtheorem{thm}{Theorem}}{}
\makeatother
\newcommand{\N}{\mathbb N}
\title[Measurable Brooks's Theorem for Directed Graphs]{Measurable Brooks's Theorem for Directed Graphs}
\author{Cecelia Higgins}
\date{Source: arXiv:2405.00991v4 (CC BY 4.0)}
\begin{document}
\maketitle

\noindent\emph{Source and licence.} Measurable Brooks's Theorem for Directed Graphs. Version arXiv:2405.00991v4 (CC BY 4.0), distributed under the Creative Commons Attribution 4.0 International licence, as recorded in the arXiv licence metadata for this exact version and shown on the abstract page https://arxiv.org/abs/2405.00991v4. Full rights notice: licenses/arxiv-2405.00991-CC-BY-4.0.txt.\par\smallskip

\noindent\textit{Editorial scope.} Counted unit: the Thm whose source label is \texttt{thm:one ended implies dicoloring} in the article (source0.tex lines 242--245, source label \texttt{thm:one ended implies dicoloring}), delivered together with the article's own complete original proof of it (source0.tex lines 247--279, one \texttt{proof} environment of 33 lines). No mathematics is written, repaired or re-derived: statement and proof are the article's own text line for line. Required context retained uncounted: thm:degplusone (lines 183--186). Every definition, display and label of those lines is retained unchanged. Omitted from this package and not redistributed: 1--182, 187--241, 246--246, 280--551. Nothing inside a retained excerpt is omitted or altered: every retained body is delivered exactly as the article's lines are written, so no omission receipt of any kind accompanies this package. Citations: the retained excerpts carry no citation command at all, so no bibliography entry of the article is transcribed and no reference is redistributed. Reference adaptation: none was needed and none was made; every reference inside the delivered unit resolves inside it, so no unresolved reference or citation is produced. Printed slips kept literal: no printed word, symbol, hypothesis or number of the counted statement or of its proof was altered. Layout and class: the article's own class is \texttt{amsart} with packages including geometry, amssymb, amsfonts, xy, enumerate, mathrsfs, tabto, tikz, tikz-cd, hyperref, verbatim, comment, appendix, graphicx; the wrapper is the lane's pinned \texttt{amsart} editorial wrapper at 10pt on A4, which restates the theorem-like environments the retained text needs and adds the article's own macro definitions that the retained text needs (\texttt{\textbackslash N}), with extra wrapper packages (bm, bbm, dsfont, xspace, cancel, mathtools). The delivered numbering is the unit's own. Assets: no retained span contains a figure environment, a graphics-inclusion command, a PDF-page inclusion, a table or a TikZ picture; no image file is copied into this package, the package declares no asset dependency and no image file is redistributed with it. Licence: version arXiv:2405.00991v4 of this article is distributed under the Creative Commons Attribution 4.0 International licence, as recorded in the arXiv licence metadata for this exact version and shown on the abstract page https://arxiv.org/abs/2405.00991v4. A full rights notice is delivered at licenses/arxiv-2405.00991-CC-BY-4.0.txt. Characters: every retained excerpt is pure ASCII, so the delivered engine, XeLaTeX with its default Latin Modern text font, reports no missing character.
\begin{thm}
\label{thm:degplusone}
    Let $D$ be a locally finite Borel digraph on a standard Borel space $X$. Let $Y$ be Polish, and let $L : X \to [Y]^{< \infty}$ be a Borel list assignment such that $\vert L(x) \vert > d^{\min}(x)$ for each $x \in X$. Then $D$ has a Borel $L$-dicoloring.
\end{thm}

\begin{thm}
\label{thm:one ended implies dicoloring}
    Let $D$ be a locally finite Borel digraph on a standard Borel space $X$, let $B \subseteq X$ be Borel, and let $f : B \rightarrow X$ be a one-ended Borel function whose graph is contained in $\tilde{D}$. Let $Y$ be Polish, and let $L : X \rightarrow [Y]^{< \infty}$ be a Borel list assignment such that $\vert L(x) \vert \geq d^{\max}(x)$ for all $x \in B$. Then $D[B]$ has a Borel $L$-dicoloring. In particular, $D[B]$ is Borel degree-list-dicolorable.
\end{thm}

\begin{proof}
    For each $n \in \N$, write $f^n[B] = \{x \in X : \text{there exist } x_1, x_2, \dots, x_n \in B \text{ such that } f(x_1) = x \text{ and } f(x_{i + 1}) = x_i \text{ for all $i < n$} \}$. Let $B_n = B \cap (f^n[B] \setminus f^{n + 1}[B])$. Note that, if $n \neq m$, then $B_n \cap B_m = \emptyset$. We claim that $B = \bigcup_{n \in \N} B_n$. Assume for contradiction that there is $x \in B \setminus \bigcup_{n \in \N} B_n$. Then $x \in f^n[B]$ for all $n \in \N$. So the set $f^{-\N}(x) = \{y \in B : \text{for some $n \geq 1$, there are } \\ x_1, x_2, \dots, x_{n - 1} \in B \text{ such that } f(x_1) = x, f(x_{i + 1}) = x_i \text{ for all $i < n - 1$, } f(y) = x_{n -1} \}$ is an infinite, finitely branching tree with root $x$. By K{\" o}nig's lemma, there is an infinite branch $(y_n)_{n \in \N}$ through $f^{-\N}(x)$. Then for each $n \in \N$, $f(y_{n + 1}) = y_n$, contradicting that $f$ is one-ended.

    As in the proof of Theorem~\ref{thm:degplusone}, we partition $X$ into countably many Borel $\tilde{D}$-independent sets $A_0, A_1, \dots$. We also fix a Borel linear ordering $<_Y$ of $Y$. Now we proceed to $L$-dicolor $D[B]$ one layer at a time. First, we work in $B_0$. For each $i \in \N$, define $B_0^i = B_0 \cap A_i$, and define $B_0^{\leq i} = \bigcup_{j \leq i} B_0^j, B_0^{< i} = \bigcup_{j < i} B_0^j$. We first color the points in $B_0^0$, then the points in $B_0^1$, and so on.

    We now define a Borel function $c_0^0 : B_0^0 \to Y$. Let $x \in B_0^0$. Then $x$ has at least one neighbor, namely $f(x)$. So since $\vert L(x) \vert \geq d^{\max}_D(x)$, it follows that $L(x) \neq \emptyset$, and so we may define $c_0^0(x)$ to be the $<_Y$-least element of $L(x)$.

    Now fix $n \in \N$, and suppose that for each $i \leq n$, a Borel function $c_0^i : B_0^{\leq i} \to Y$ has been defined so that:
    \begin{itemize}
        \item If $j < i$, then $c_0^i \restriction B_0^{\leq j} = c_0^j$;
        \item For all $x \in B_0^{\leq i}$, $c_0^i(x) \in L(x)$; and
        \item For all $x \in B_0^{\leq i}$, if $x$ has an out-neighbor $y \in B_0^{< i}$ and an in-neighbor $z \in B_0^{< i}$ with $\alpha = c_0^{i - 1}(y) = c_0^{i - 1}(z)$, then $c_0^i(x) \neq \alpha$.
    \end{itemize}
    Now for each $x \in B_0^{n + 1}$, let
    $$L_0^{n + 1}(x) = L(x) \setminus \{\alpha \in Y : \exists y, z \in B_0^{\leq n} \text{ such that } y \in N^+(x), z \in N^-(x), \text{and } c_0^n(y) = c_0^n(z) = \alpha \}.$$
    Then since $x$ has at least one neighbor, namely $f(x)$, not contained in $B_0$, we have that $\{\alpha \in Y : \exists y, z \in B_0^{\leq n} \text{ such that } y \in N^+(x), z \in N^-(x), \text{and } c_0^n(y) = c_0^n(z) = \alpha \}$ has cardinality at most $d^{\max}_D(x) - 1$, and so
    \begin{align*}
        \vert L_0^{n + 1}(x) \vert & \geq \vert L(x) \vert - (d^{\max}_D(x) - 1) \\
        & \geq d^{\max}_D(x) - (d^{\max}_D(x) - 1) \\
        & = 1.
    \end{align*}
    Therefore, we may extend $c_0^n : B_0^{\leq n} \to Y$ to a function $c_0^{n + 1} : B_0^{\leq (n + 1)} \to Y$ by letting $c_0^{n + 1}(x)$ be the $<_Y$-least element of $L_{n + 1}(x)$ for each $x \in B_{n + 1}$. After this recursive procedure is complete, we define $c_0 = \bigcup_{n \in \N} c_0^n$.

    We now work on $B_1$. As before, we partition $B_1$ into independent sets $B_1^0, B_1^1, \dots$. To define $c_1^0 : B_1^0 \to Y$, we first re-initialize the list assignments by setting
    $$L_1^0(x) = L(x) \setminus \{\alpha \in Y : \exists y, z \in B_0 \text{ such that } y \in N^+(x), z \in N^-(x), \text{and } c_0(y) = c_0(z) = \alpha \}$$
    for each $x \in B_1^0$. Then as before, since for each $x \in B_1^0$ there is a neighbor $f(x)$ of $x$ that does not belong to $B_1$, we have $\vert L_1^0(x) \vert \geq 1$, and so we may define $c_1^0(x)$ to be the $<_Y$-least element of $L_1^0(x)$. Given $n \in \N$, having constructed $c_1^0, \dots, c_1^n$, we define $c_1^{n + 1}$ on $B_1^{n + 1}$ as follows: For a fixed $x \in B_1^{n + 1}$, set
    \begin{align*}
        L_1^{n + 1}(x) = L(x) \setminus \{\alpha \in Y : \exists y, z \in (B_0 \cup B_1^{\leq n}) \text{ such that } y \in N^+(x), z \in N^-(x), \\ \text{and } (c_0 \cup c_1^n)(y) = (c_0 \cup c_1^n)(z) = \alpha \}.
    \end{align*}
    Then since $f(x) \notin B_1$, we have $\vert L_1^{n + 1}(x) \vert \geq 1$. Thus, we may define $c_1^{n + 1}(x)$ to be the $<_Y$-least element of $L_1^{n + 1}(x)$. Once this procedure is complete, define $c_1 = \bigcup_{n \in \N} c_1^n$. Continue on in this way to the sets $B_2, B_3, \dots$, constructing functions $c_2, c_3, \dots$.

    Finally, let $c = \bigcup_{n \in \N} c_n$. Then $c : B \to Y$ is Borel. Assume for contradiction that there is a $c$-monochromatic dicycle $C = (x_0, x_1, \dots, x_k = x_0)$. Let $n \in \N$ be maximal with $C \cap B_n \neq \emptyset$, and let $m \in \N$ be maximal with $C \cap B_n^m \neq \emptyset$. Then there is $i < k$ with $x_i \in B_n^m$. It follows that $x_{i - 1}, x_{i + 1} \notin B_n^m$. So either $x_{i - 1} \in B_{n'}$ for some $n' < n$, or $x_{i - 1} \in B_n^{m'}$ for some $m' < m$, and similarly for $x_{i + 1}$. In any case, it follows that $\alpha = c(x_{i - 1}) = c(x_{i + 1})$ is not an element of $L_n^m(x_i)$, a contradiction since $c(x_i) = \alpha$. Thus, $c$ is a Borel $L$-dicoloring of $D[B]$.
\end{proof}

\end{document}
